\documentclass[12pt]{article}
\usepackage{amsmath,amssymb,amsthm,geometry,hyperref}
\geometry{a4paper,margin=1in}
\newtheorem{theorem}{Theorem}
\newtheorem{lemma}[theorem]{Lemma}
\title{Project Euler Problem 482: The Incenter of a Triangle}
\author{}
\date{}
\begin{document}
\maketitle

\section{Problem Statement}
$ABC$ is an integer-sided triangle with incenter $I$ and perimeter $p$. The segments $IA$, $IB$, $IC$ have integral length. Let $L = p + |IA| + |IB| + |IC|$ and $S(P) = \sum L$ for all such triangles with $p \le P$. Given $S(10^3) = 3619$, find $S(10^7)$.

\section{Mathematical Foundation}

\begin{theorem}[Incenter distance formula]
For a triangle with sides $a,b,c$, semiperimeter $s=(a+b+c)/2$, and incenter $I$:
\[
|IA|^2 = \frac{bc(s-a)}{s},\quad |IB|^2 = \frac{ac(s-b)}{s},\quad |IC|^2 = \frac{ab(s-c)}{s}.
\]
\end{theorem}

\begin{proof}
The inradius is $r = K/s$ where $K = \sqrt{s(s-a)(s-b)(s-c)}$. Using $|IA| = r/\sin(A/2)$ and the half-angle identity $\sin(A/2) = \sqrt{(s-b)(s-c)/(bc)}$:
\[
|IA| = \frac{K}{s}\cdot\sqrt{\frac{bc}{(s-b)(s-c)}} = \sqrt{\frac{(s-a)\cdot bc}{s}}.
\]
The other two follow by cyclic symmetry.
\end{proof}

\begin{lemma}[Change of variables]
Setting $x = s-a$, $y = s-b$, $z = s-c$ (so $a=y+z$, $b=x+z$, $c=x+y$, $s=x+y+z$):
\begin{align}
|IA|^2 &= \frac{x(x+y)(x+z)}{x+y+z},\\[4pt]
|IB|^2 &= \frac{y(x+z)(y+z)}{x+y+z},\\[4pt]
|IC|^2 &= \frac{z(x+y)(y+z)}{x+y+z}.
\end{align}
Each must be a perfect square for the incenter distances to be integers.
\end{lemma}

\begin{proof}
Direct substitution into Theorem~1.
\end{proof}

\begin{lemma}[Parity constraint]
The perimeter $p = 2(x+y+z)$ is even, and $x,y,z \ge 1$ automatically satisfy the triangle inequality.
\end{lemma}

\begin{proof}
For $x,y,z$ to be positive integers, $s$ must be a positive integer, so $p=2s$ is even. The inequality $a < b+c$ is equivalent to $x > 0$, which holds.
\end{proof}

\begin{theorem}[Product identity]
\[
|IA|^2\cdot|IB|^2\cdot|IC|^2 = \frac{xyz\,(x+y)^2(y+z)^2(x+z)^2}{(x+y+z)^3}.
\]
\end{theorem}

\begin{proof}
Multiply the three expressions from Lemma~1.
\end{proof}

\section{Algorithm}
\begin{enumerate}
\item Enumerate positive integers $x \le y \le z$ with $2(x+y+z) \le P$.
\item For each triple, verify $s \mid x(x+y)(x+z)$ and that the quotient is a perfect square; similarly for $|IB|^2$ and $|IC|^2$.
\item For valid triples, compute $L = 2s + |IA| + |IB| + |IC|$ and account for distinct permutations.
\item Accumulate $S(P)$.
\end{enumerate}

\section{Complexity Analysis}
\begin{itemize}
\item \textbf{Time:} $O(P^{3/2})$ worst case, heavily pruned by divisibility and perfect-square tests.
\item \textbf{Space:} $O(1)$ auxiliary.
\end{itemize}

\section{Answer}

\[
\boxed{1400824879147}
\]

\end{document}
